\documentclass[../../main.tex]{subfiles}

\begin{document}

\section{Roots of Unity}






\begin{problem}{1976 USAMO Problem 5 \medium}{1976 USAMO Problem 5}
For polynomials $P(x)$, $Q(x)$, $R(x)$, and $S(x)$ that satisfy
the following equation,
\[
P(x^5) + xQ(x^5) + x^2 R(x^5)
= (x^4 + x^3 + x^2 + x + 1) S(x),
\]
show that $x - 1$ is a factor of $P(x)$.
\end{problem}
\begin{solution}
Looking at the equation, we can see $x^4 + x^3 + x^2 + x + 1$ and
$x^5$. This should remind the equation $x^5 - 1 = (x - 1)
(x^4 + x^3 + x^2 + x + 1)$. In other words, we should be reminded of
the fifth roots of unity where we define $\omega_5 = 1$ and
$\omega_1 + \omega_2 + \cdots + \omega_5 = 0$.

The next thing that we want to try is substituting $\omega_1, \cdots,
\omega_4$ since we know that the equation always hold.
\begin{align*}
    P(1) + \omega_1 Q(1) + \omega_1^2 R(1) &= 0 \\
    P(1) + \omega_2 Q(1) + \omega_2^2 R(1) &= 0 \\
    P(1) + \omega_3 Q(1) + \omega_3^2 R(1) &= 0 \\
    P(1) + \omega_4 Q(1) + \omega_4^2 R(1) &= 0
\end{align*}
Now we see that only $P(1)$ has equal coefficient across all four
equations. We also have defined $\omega_1, \cdots \omega_5$ to be
the fifth root of unity. Meaning, if we multiply the equations
appropriately, we would be able to cancel out $Q(1)$ and $R(1)$.
\begin{align*}
    \omega_1 P(1) + \omega_1^2 Q(1) + \omega_1^3 R(1) &= 0 \\
    \omega_2 P(1) + \omega_2^2 Q(1) + \omega_2^3 R(1) &= 0 \\
    \omega_3 P(1) + \omega_3^2 Q(1) + \omega_3^3 R(1) &= 0 \\
    \omega_4 P(1) + \omega_4^2 Q(1) + \omega_4^3 R(1) &= 0
\end{align*}
Subtracting the four equations from the original equations,
we obtain the following equations.
\begin{align*}
    &\qquad\quad P(1) + \omega_1 Q(1) + \omega_1^2 R(1) \\
    &-\left(
        \omega_1 P(1) + \omega_1^2 Q(1) + \omega_1^3 R(1)
    \right) \\
    &+ \quad\ \ \ P(1) + \omega_2 Q(1) + \omega_2^2 R(1) \\
    &-\left(
        \omega_2 P(1) + \omega_2^2 Q(1) + \omega_2^3 R(1)
    \right) \\
    &+ \quad\ \ \ P(1) + \omega_3 Q(1) + \omega_3^2 R(1) \\
    &-\left(
        \omega_3 P(1) + \omega_3^2 Q(1) + \omega_3^3 R(1)
    \right) \\
    &+ \quad\ \ \ P(1) + \omega_4 Q(1) + \omega_4^2 R(1) \\
    &-\left(
        \omega_4 P(1) + \omega_4^2 Q(1) + \omega_4^3 R(1)
    \right) = 0
\end{align*}
Solving the equation, we get $5 P(1) = 0$ and $x - 1$ is a factor
of $P(x)$.
\end{solution}






\subsection{Roots of Unity Filter}

\[
a_0 + a_n + a_{2n} + \cdots
= \frac{f(\omega) + f\left(\omega^2\right) + \cdots +
    f\left(\omega^n\right)
}{n}
\]

\begin{problem}{}{}
In Enoch's dreamland, there are $999$ mysterious boxes each labeled
uniquely from $1$ to $999$. If Enoch chooses the box(es) randomly
and the sum of the number of chosen box(es) is divisible by $3$,
then he earns a dollar from the owner of the boxes. However, he is
not granted money if he chooses the same combination next time.
What is the maximum amount of money Enoch can earn?
\end{problem}
\begin{solution}
This problem is basically asking for the number of ways to choose
numbers randomly from $1$ through $999$ such that their sum is
divisible by $3$. In other words, how many subsets of the set
$\{1, 2, \ldots, 999\}$ have the sum of their elements divisible by
three?

When a large sequence appears, it is reasonable to generate functions.
For the function that we are generating, we will associate the subsets
with the its coefficients. This allows an algebraic approach towards
the combinatorial problem.

Consider the following function.
\begin{align*}
    f(x)
    &= \prod_{i=1}^{999} \left( 1 + x^i \right)
    = (1 + x) \left( 1 + x^2 \right) \left( 1 + x^3 \right) \cdots
        \left( 1 + x^{999} \right)
\end{align*}
Notice that each number in the subset corresponds to the exponents of
$x$. Consider the following samples to see how they are related.
\begin{align*}
    &\text{Exponent: $3$} \Rightarrow \text{Coefficient: $2$} \\
    &\left.\begin{aligned}
        3 &= 3 &&\Rightarrow \{ 3 \} \\
        3 &= 1 + 2 &&\Rightarrow \{ 1, 2 \}
    \end{aligned}\right\}
    \quad \text{2 Cases}
\end{align*}
\begin{align*}
    &\text{Exponent: $6$} \Rightarrow \text{Coefficient: $4$} \\
    &\left.\begin{aligned}
        6 &= 6 &&\Rightarrow \{ 6 \} \\
        6 &= 1 + 5 &&\Rightarrow \{ 1, 5 \} \\
        6 &= 2 + 4 &&\Rightarrow \{ 2, 4 \} \\
        6 &= 1 + 2 + 3 &&\Rightarrow \{ 1, 2, 3 \} \\
    \end{aligned}\right\}
    \quad \text{4 Cases}
\end{align*}
The process of investigating the subsets with sum $n$ perfectly
mirrors the behavior of terms of $x^n$ and its coefficients.
Therefore, the sum of coefficients of $x^{3k}$ where
$k \in \mathbb{N}_0$ is the number of subsets whose sum of the
elements are divisible by $3$.

Continuing, we can rewrite the function $f$.
\[
f(x) = a_0 + a_1x^1 + a_2x^2 + \cdots + a_{499500}x^{499500}
\]
Using the third roots of unity with $\omega_1, \omega_2, \omega_3$,
\begin{align*}
    f(\omega_1)
    &= a_0 + a_1\omega_1 + \cdots + a_{499500}\omega_1^{499500} \\
    f(\omega_2)
    &= a_0 + a_1\omega_2 + \cdots + a_{499500}\omega_2^{499500} \\
    f(\omega_3)
    &= a_0 + a_1\omega_3 + \cdots + a_{499500}\omega_3^{499500}
    \\[0.5em]
    f(\omega_1) + f(\omega_2) + f(\omega_3)
    &= 3a_0 + 3a_1(\omega_1 + \omega_2 + \omega_3) + \cdots \\
    &\qquad + \, a_{499500} \left(
        \omega_1^{499500} + \omega_2^{499500} + \omega_3^{499500}
    \right) \\
    &= 3a_0 + 3a_3 + 3a_6 + \cdots + 3a_{499500}
\end{align*}
holds since $e^{ i \frac{0\pi}{3} }, e^{ i \frac{2\pi}{3} },
e^{ i \frac{4\pi}{3} }$ will continue to repeat. Therefore, it
suffices to find the values of $\frac{f(\omega_1) + f(\omega_2) +
f(\omega_3)}{3}$. Using the period again, we can obtain the following.
\begin{align*}
    \frac{ f(\omega_1) + f(\omega_2) + f(\omega_3) }{3}
    &= \frac{ \left(
        \left( 1 + \omega_1 \right)
        \left( 1 + \omega_1^2 \right)
        \left( 1 + \omega_1^3 \right)
    \right)^{333}}{3} \\
    &\qquad +\, \frac{ \left(
        \left( 1 + \omega_2 \right)
        \left( 1 + \omega_2^2 \right)
        \left( 1 + \omega_2^3 \right)
    \right)^{333}}{3} \\
    &\qquad +\, \frac{ \left(
        \left( 1 + \omega_3 \right)
        \left( 1 + \omega_3^2 \right)
        \left( 1 + \omega_3^3 \right)
    \right)^{333}}{3}
\end{align*}
Without loss of generality, assume that $\omega_1 =
e^{ i \frac{0\pi}{3} }$, $\omega_2 = e^{ i \frac{2\pi}{3} }$, and
$\omega_3 = e^{ i \frac{4\pi}{3} }$. Except the term
$\left( 1 + \omega_1 \right) \left( 1 + \omega_1^2 \right)
\left( 1 + \omega_1^3 \right)$, the other expressions
could be rewritten as $(1 + \omega_1)(1 + \omega_2)(1 + \omega_3)$.
Moreover since $1 + \omega_i$ is the root of $(x - 1)^3 = 1$,
\[
(1 + \omega_1)(1 + \omega_2)(1 + \omega_3) = 2.
\]
by Vieta's formula. Continuing,
\[
\frac{f(\omega_1) + f(\omega_2) + f(\omega_3)}{3}
= \frac{2^{999} + 2^{333} + 2^{333}}{3}
= \frac{2^{334} + 2^{999}}{3}
\]
Therefore, Enoch can earn $\frac{2^{334} + 2^{999}}{3}$ dollars in
total.
\end{solution}






\end{document}


